% ===== §9 =====
\section{East and West: An Unactivated Resonance}
\label{p24:sec:eastwest}

\noindent
Before the antithesis states its framework, one connection across the traditions just surveyed must be drawn, because it is the connection that motivates the framework and shows why a new one is needed. The six traditions do not merely occupy different coordinates on the map; two of them, arising on opposite sides of the world and out of wholly unrelated lineages, point toward the same buried truth, and the recognition of this convergence is the hinge on which the paper turns from thesis to antithesis. The section that performs this turn is not a comparison offered for its own sake, a pairing of East and West laid side by side for balance. It is the identification of a resonance that has never been activated, and the claim that activating it requires precisely the apparatus the next section will supply.

\subsection*{The resonance}

Set the dialogic tradition beside the verbal force of \emph{hua}, and a single shape appears twice. Bakhtin's dialogism holds that meaning is generated in the between of address and answer and is constitutively unfinished, that nothing conclusive has yet been said and the ultimate word cannot be spoken \parencite{bakhtin1984poetics}. The Chinese \emph{hua} holds, in the plain grammar of the term, that at the heart of culture stands a verb of transformation, an ongoing bringing-to-form rather than a settled deposit. These are not the same doctrine, and it would be a distortion to collapse them; the one is a twentieth-century poetics of the novel and of the utterance, the other a classical vocabulary of civilization and cultivation. But they converge on one thing, and they converge on it against the main current of the traditions around them: that culture is \emph{process before it is product}, generation before deposit, the unfinished before the finished. What the settled traditions of the West most persistently lose---the moment of generation, the ongoing act, the refusal of closure---these two, from their opposite origins, independently preserve. The resonance is real, and it is the more striking for being unforced: a Russian theory of the dialogic word and a classical Chinese verb of formation, with no line of influence between them, agree that the being of culture is generative.

\subsection*{Why the resonance has not been activated}

Yet the resonance has never done the work it could do, and the reason is that each of its two poles possesses the insight without the means to develop it. The Chinese tradition holds the generative conception of culture in a vocabulary of great depth, but it does not bring to that conception the analytic apparatus the modern question of culture requires: it does not thematize drive and its lack, it does not articulate the registers of experience, and it does not carry a critique of the material and political conditions under which culture is produced and captured. It has the insight and not the analysis. The Western traditions, conversely, have built exactly these analytic instruments---the psychoanalysis of desire, the theory of the symbolic order, the political economy of cultural production---but have overwhelmingly put them at the service of the settled conception, using them to anatomize culture as a prior web, a deep structure, or a congealed order of power, and thereby burying the very generativity the dialogic strand alone, at the margin, kept alive. The West has the analysis and, mostly, not the insight. Each pole, in short, has half of what a complete account requires, and the two halves have never been joined, because joining them requires a framework that is neither simply Chinese nor simply Western, that can hold the generative conception of culture and the analytic apparatus of drive, register, and material condition in a single structure.

\subsection*{The framework as the activator}

This is the office the next section's framework is built to perform, and naming that office is the business of this turn. The three-register dialectic of the relational Real is not offered as a reconciliation that splits the difference between East and West, nor as a synthesis that dissolves either into the other. It is offered as the apparatus that activates the resonance: a more fundamental ontology of relational generation within which each tradition's truth finds its place. The generative conception that the dialogic strand and the Chinese \emph{hua} independently preserve becomes, within this framework, the basic character of culture as such---culture as the ongoing, unfinished exteriorization of a relational generativity. And the analytic instruments the West has built become the means of specifying \emph{how} that generation works: drive as what is coded and sublimated in it, the three registers as the dimensions across which it is exteriorized, the material conditions as what constrains it, the phallic organization of the symbolic as where power enters it, and the counter-power of a just generation as what keeps that power from entrenching. The framework does not adjudicate between the traditions; it gives each its function within a generative ontology none of them, alone, possessed. With this office named, the thesis is complete and the antithesis can begin, not by adding a seventh view to the six, but by stating the framework within which all six are to be reread and each assigned its truth and its limit.
